% Template for scripts/build_codebook.py. Double-angle-bracket markers are filled from the corpus
% CSV on every build; everything else is fixed text. Edit wording here, not in a generated .tex.
\documentclass[10pt]{article}
\usepackage[T1]{fontenc}
\usepackage[utf8]{inputenc}
\usepackage{lmodern}
\usepackage[margin=0.9in]{geometry}
\usepackage{booktabs,longtable,array,tabularx}
\usepackage[dvipsnames]{xcolor}
\usepackage[hidelinks]{hyperref}

\let\olditemize\itemize
\renewcommand{\itemize}{\olditemize\setlength{\itemsep}{1pt}\setlength{\parskip}{0pt}}
\setlength{\parindent}{0pt}
\setlength{\parskip}{0.55em}
\renewcommand{\arraystretch}{1.18}
\newcommand{\col}[1]{\texttt{#1}}
\newcolumntype{L}[1]{>{\raggedright\arraybackslash}p{#1}}

\pagestyle{myheadings}
\markright{\small U.S. Secretary of State Speech Corpus --- Codebook, release <<VERSION>>}

\begin{document}

\begin{center}
{\LARGE\bfseries U.S. Secretary of State Speech Corpus}\\[0.4em]
{\Large Codebook for the dataset columns}\\[0.6em]
{\normalsize Reference release \textbf{<<VERSION>>} \quad$\cdot$\quad <<N_ROWS>> rows \quad$\cdot$\quad <<N_COLS>> columns \quad$\cdot$\quad <<DATE_MIN>> to <<DATE_MAX>>}
\end{center}

\section*{1\quad The dataset in brief}

Each row is one public statement by a U.S.\ Secretary of State: a speech,
remarks, interview, press availability, testimony, press statement or op-ed,
collected from \texttt{state.gov} and its archived sites
(\texttt{1997-2001.state.gov} through \texttt{2021-2025.state.gov}).
The corpus covers <<SEC_FIRST>> through <<SEC_LAST>>. New Rubio material is added by a weekly
automated update, and each update is published as a new versioned release.

There are 19 columns. Ten \emph{base} columns describe the document and hold its
text (Section~3). Nine \emph{derived} columns are computed from the base columns by a
separate enrichment stage (Section~4). The enrichment stage never changes a base
column; the release verifier checks this (Check~Q).

\paragraph{File format.} One UTF-8 CSV file with a header row and standard
double-quote escaping (RFC~4180). The text fields contain line breaks inside
quoted values, so read the file with a real CSV parser, not line by line. Some
fields are very long (the largest \col{Raw Text} is over <<RAW_MAX_CHARS>> characters), so
raise the parser's field-size limit:
\begin{itemize}
\item Python \texttt{csv}: call \texttt{csv.field\_size\_limit(10**8)}.
\item pandas: \texttt{pd.read\_csv(path, dtype=str, keep\_default\_na=False)} keeps
empty list columns as empty strings rather than \texttt{NaN}.
\item R: \texttt{readr::read\_csv(path)} handles embedded newlines.
\end{itemize}
Each release ships with a sidecar file, \texttt{<file>.manifest.json}. It records
the row count, the column order, the SHA-256 of the CSV, and a fingerprint of every
row's \col{Row Hash} (Section~5).

\paragraph{List-valued columns.} Five derived columns hold lists. Items are
joined by a semicolon and a space (\texttt{"; "}), keep their order of first
appearance, and contain no duplicates within a row. An empty string means the
list is empty. It does not mean the value is missing.

\paragraph{Which text column to use.}
\begin{itemize}
\item \col{SOS Text}: the Secretary's own words only. Use it for analyses of what the Secretary said.
\item \col{Natural Text}: the whole cleaned document, with every speaker. Use it for context, or for the questions the Secretary was asked.
\item \col{Raw Text}: the page as scraped. Use it only for provenance and auditing; it still contains site navigation and encoding errors.
\end{itemize}

\section*{2\quad Coverage}

\begin{center}
\begin{tabular}{lrll}
\toprule
\col{Secretary} & Rows & First date & Last date \\
\midrule
<<COVERAGE_ROWS>>
\bottomrule
\end{tabular}
\end{center}
{\small Tenures overlap by a few days because a Secretary-designate's confirmation
testimony is dated before they were sworn in (e.g.\ <<OVERLAP_EXAMPLE>>). \col{Secretary} identifies who spoke, not who held the office on
that date.}

\section*{3\quad Base columns}

\begin{longtable}{@{}L{1.25in}L{0.75in}L{4.35in}@{}}
\toprule
Column & Type & Definition and coding rules \\
\midrule
\endhead
\col{Row Id} & integer & Row position, 0 to $N-1$, with no gaps. \textbf{Renumbered in every
release}, so the same speech can have a different \col{Row Id} in another release.
Use \col{URL} or \col{Row Hash} as a stable key across releases. \\
\col{Secretary} & categorical & The speaking Secretary's last name: <<SEC_LIST>>. \\
\col{Date} & date & \texttt{YYYY-MM-DD} (ISO 8601). The date of the statement as given on the page.
When the page date and the URL's \texttt{/YYYY/MM/} path disagree and the page text gives a
date that matches the URL month, that date is used. The verifier checks that no fixable
mismatch remains. \\
\col{Title} & string & The page title. The site suffix and any date prefix are removed,
whitespace is normalised, and the first letter is capitalised. Not unique, because many
statements share a generic title. \\
\col{Speech Type} & categorical & One of <<N_TYPES>> controlled values. Row counts are in the table
after this one. \\
\col{Location} & string & Exactly one comma: \texttt{City, Region}. Region is a country, or a
U.S.\ state for U.S.\ cities (e.g.\ \texttt{Beijing, China}; \texttt{New York City, New York}).
The value is \texttt{Washington, D.C.} when the page itself gives a Washington newsroom origin
(``Office of the Spokesperson/Spokesman'', ``Via Teleconference''), and also for
statements given in Washington. A location the pipeline cannot resolve is \emph{not} given
a default. The row is held back (quarantined) for manual review instead of being published.
\texttt{Washington, D.C.} is <<WASH_PCT>>\% of rows. There are no ``En Route'' or non-geographic values. \\
\col{URL} & string & The source URL exactly as first scraped. \textbf{Fixed once a row
is released}: it is an input to \col{Row Hash}, so it is never rewritten even after the page
moves. Use it to cite and trace a row's provenance. Some older pages may now be
available only from the Internet Archive. \\
\col{SOS Text} & text & Only the Secretary's own words. The text of other speakers
(questioners, moderators, foreign ministers) is removed, as are speaker labels, site navigation and
footers, ``Released on \ldots'' date stamps, document reference codes, audience-reaction notes
such as \texttt{(Laughter.)} and \texttt{(Applause.)}, and \texttt{(inaudible)} markers.
Encoding errors are repaired. For single-speaker documents (most press statements, many
remarks), it is identical to \col{Natural Text}; this is true for <<SOS_EQ_NAT>> rows (<<SOS_EQ_NAT_PCT>>\%). \\
\col{Natural Text} & text & The full cleaned document with \emph{all} speakers, in order.
Speaker labels are \textbf{kept} (e.g.\ \texttt{SECRETARY KERRY:}, \texttt{QUESTION:}),
so it can be split into turns. Site navigation and footer text are removed, and encoding
errors are repaired. \\
\col{Raw Text} & text & The page text as scraped, with no cleaning. It includes site
navigation, archive banners and footers, and any encoding errors (e.g.\ the
\texttt{\~{a}\textemdash\textemdash} separator that is removed from the two cleaned text fields). For a small set of
rows whose original page was broken, this is the text of a re-fetched copy (the current page or
an Internet Archive snapshot). Use it for provenance and auditing, not for analysis. \\
\bottomrule
\end{longtable}

\begin{center}
\small
\begin{tabular}{lr@{\qquad}lr@{\qquad}lr}
\toprule
\col{Speech Type} & Rows & \col{Speech Type} & Rows & \col{Speech Type} & Rows \\
\midrule
<<SPEECH_TYPE_ROWS>>
\bottomrule
\end{tabular}
\end{center}

\section*{4\quad Derived columns}

These columns are added by the enrichment stage (\texttt{corpus\_enrich\_v1.py}).
The last four columns before \col{Canonical URL} come from automatic named-entity
recognition and \textbf{have not been reviewed by hand}. Treat them as approximate
indexes for finding and filtering documents, not as validated measures.

\begin{longtable}{@{}L{1.25in}L{0.75in}L{4.35in}@{}}
\toprule
Column & Type & Definition and coding rules \\
\midrule
\endhead
\col{SOS Word Count} & integer & Number of whitespace-separated tokens in \col{SOS Text}
(Python \texttt{len(text.split())}). Mean <<SOS_WC_MEAN>>, median <<SOS_WC_MEDIAN>>, range <<SOS_WC_MIN>>--<<SOS_WC_MAX>>. <<SOS_WC_SHORT>> rows have
fewer than 20 words; these are flagged by the verifier but kept. \\
\col{Natural Word Count} & integer & The same count for \col{Natural Text}. Mean <<NAT_WC_MEAN>>,
median <<NAT_WC_MEDIAN>>. It is normally $\geq$ \col{SOS Word Count}. \\
\col{Row Hash} & string & 64-character lowercase hex SHA-256 of nine fields:
\col{Secretary}, \col{Date}, \col{Title}, \col{Speech Type}, \col{Location}, \col{URL},
\col{SOS Text}, \col{Natural Text}, \col{Raw Text}. The fields are joined in that order by
the unit separator \texttt{U+001F} and encoded as UTF-8. \col{Row Id} is excluded, so the
hash stays the same when rows are renumbered. It changes if any listed field changes. Unique
within a release. \\
\col{Speakers Present} & list & Every distinct speaker in \col{Natural Text}, in order of
first appearance. Never empty: the row's own Secretary is always included, by last name.
Unnamed roles are shown as bracketed placeholders (<<SENTINELS>>, \ldots) so they are not mistaken for names. Named speakers appear as they are
labelled in the transcript (usually by last name). Mean <<SPEAKERS_MEAN>> speakers per row. \\
\col{People Mentioned (Natural Text)} & list & People \emph{referred to} anywhere in
\col{Natural Text}, found by spaCy \texttt{en\_core\_web\_lg} PERSON entities. Names used
only to address someone directly (``Thank you, Mr.\ Chairman'') are excluded by a set of
rules. Names are as written: the same person can appear as a full name and as a surname
(e.g.\ \texttt{Carl Bildt; Carl; Bildt}). Empty for <<PEOPLE_NAT_EMPTY>> rows. \\
\col{People Mentioned (SOS Text)} & list & The subset of the previous column whose
mentions fall inside the Secretary's own turns. This is always a subset of
\col{People Mentioned (Natural Text)}. Empty for <<PEOPLE_SOS_EMPTY>> rows. \\
\col{Countries Mentioned} & list & Countries named in \col{Natural Text}. spaCy GPE/LOC
entities are mapped to a standard country name using a curated alias table, followed by an
\emph{exact} match against ISO 3166 (no fuzzy matching). U.S.\ states and territories are
mapped to \texttt{United States}. Some territories are listed separately (e.g.\ \texttt{Hong
Kong}). Historical states are kept (\texttt{Soviet Union}, \texttt{Yugoslavia},
\texttt{Czechoslovakia}). Mentions that cannot be resolved are dropped. Empty for <<COUNTRIES_EMPTY>> rows. \\
\col{Country ISO Codes Mentioned} & list & ISO 3166-1 alpha-3 codes, \textbf{in the same
positions} as \col{Countries Mentioned}: item $i$ is the code for country $i$. Where a
country has no ISO code (<<NO_ISO_LIST>>), the code is
\texttt{N/A}, so the two lists always have the same length. \\
\col{Canonical URL} & string & Where the page is \emph{now}. It equals \col{URL} unless
\texttt{state.gov} has since moved the page. In September 2026, for example,
\texttt{/releases/office-of-the-spokesperson/} became \texttt{/releases/office-of-the-spokesman/};
<<CANON_DIFF>> rows differ for this reason. It is not part of the hash, so it can be updated when the
site changes again. Cite \col{URL}; use \col{Canonical URL} to fetch the live page. \\
\bottomrule
\end{longtable}

\noindent\begin{minipage}{\linewidth}
\paragraph{Example row} (Row Id <<EX_ROW_ID>> in <<VERSION>>; text fields shortened):
\begin{center}
\small
\begin{tabularx}{\linewidth}{@{}L{1.9in}X@{}}
\toprule
<<EXAMPLE_ROWS>>
\bottomrule
\end{tabularx}
\end{center}
\end{minipage}

\section*{5\quad Integrity and verification}

\begin{itemize}
\item \textbf{Manifest.} \texttt{<file>.manifest.json} lists the row count, the column
order, the SHA-256 of the CSV file (\texttt{csv\_sha256}) and a fingerprint of the whole
dataset (a hash of every \col{Row Hash} in row order). For <<VERSION>> the file hash begins
\texttt{<<SHA_HEAD>>} and ends \texttt{<<SHA_TAIL>>}. Compare it with \texttt{sha256sum} to confirm
you have the exact file.
\item \textbf{Row-level edits.} Recompute \col{Row Hash} with the formula in Section~4 to
check that a row has not been changed since release.
\item \textbf{Release verifier.} Every release must pass
\texttt{scripts/corpus\_verify\_research\_grade.py} with no failures before it is published.
The verifier checks the schema and controlled vocabularies, duplicates (by URL, by URL slug and
date, and by content), leftover site text and encoding errors in the cleaned text fields,
the \col{Location} format, that word counts and hashes match when recomputed, and the
format and internal consistency of every derived column.
\item \textbf{Patches.} A text fix applied after release changes \col{SOS Text} and
\col{Natural Text} only. Each fix is listed in a patch audit file
(\texttt{corpus\_v<version>\_patch\_audit.csv}) that is published with the release.
\end{itemize}

\end{document}
